\documentclass[12pt]{article}
% KTU template's existing natbib / unsrtnat configuration is provided by this style.
% Build from repository root; this project does not alter root main.tex or variables.tex.
\usepackage{ktu_thesis_style}
\usepackage{float}
\usepackage{microtype}
\usepackage{amssymb}
\usepackage{xcolor}
\usepackage{pdflscape}
\usepackage{array}
\usepackage{etoolbox}
\hypersetup{hidelinks,pdftitle={UV PCB Conformal Coating Inspection Literature Review and Methodology}}
\setcounter{secnumdepth}{3}
\setcounter{tocdepth}{2}
\newcommand{\ccl}{conformal coating}
\newcommand{\systemname}{UV-CCI}
\newcommand{\evidencenote}[1]{\textit{Evidence caveat: #1}}
\begin{document}
\begin{center}
{\large\bfseries KAUNAS UNIVERSITY OF TECHNOLOGY}\par\vspace{0.6cm}
{\large\bfseries Scientific Literature Review and Experimental Methodology}\par\vspace{1.0cm}
{\Large\bfseries Development and Evaluation of a Low-Cost Reference-Aware UV Fluorescence Inspection Architecture for PCB Conformal Coating}\par\vspace{0.9cm}
{\normalsize Research working manuscript --- version 1.0, 23 October 2026}\par\vspace{0.5cm}
{\small Scope: nominal 70 $\times$ 70 mm PCB; academic research and possible industrial validation}\par
\end{center}
\vspace{1.0cm}
\noindent\textbf{Document status.} Evidence-grounded literature synthesis and proposed methodology. This manuscript does \emph{not} present acquired UV coating images, trained models or independently measured inspection results. Bibliographic identities were checked through public publisher/standard/vendor/dataset records; comprehensive systematic full-text appraisal remains in progress. Numeric claims from manufacturer brochures are labelled accordingly.\par
\vspace{0.3cm}
\noindent\textbf{Data confidentiality.} No employer source code, proprietary PCB images or production process parameters are included. Any eventual publication requires review of intellectual property and information-release permissions.\par
\clearpage
\begin{abstract}
Inspection of transparent conformal coating is a physically constrained machine-vision problem: UV excitation can reveal fluorescent coating, but visible brightness is also affected by excitation irradiance, camera response, optical filtering, material composition, geometry and shadows. Research already demonstrates ultraviolet bubble detection with CNN and transfer learning, automated UV illumination suitability selection, and CNN-based inspection of dewetting and capillary defects. Commercial coating inspection systems combine UV lighting, reference images, CAD coating plans, additional camera angles and optional independent thickness metrology. Thus, applying a neural network to UV-coated PCBs is not in itself novel. This study develops a traceable synthesis of directly relevant evidence, transferable computer-vision methods and standards, then proposes an evaluable low-cost architecture for fixed 70~mm square printed circuit board assemblies. It combines radiometric correction, geometry registration, a probabilistic coating representation, a manufacturing reference map, morphological defect characterization and risk-sensitive PCB-level decisions. A prospective experimental design compares classical, standalone learned and hybrid approaches on physical-board- and lot-disjoint test groups, with explicit independent ground truth, defect-specific sensitivity, false acceptance, uncertainty, robustness and whole-cycle CPU timing. The report also identifies the gap between fluorescence coverage and quantitative coating thickness, the restricted relevance of popular PCB defect datasets, and the limits of current source access. The output is an experimentally actionable research protocol and preliminary academic literature review, not a validated production inspection system.
\end{abstract}
\tableofcontents
\clearpage
\section{Research problem, objectives and boundaries}
\subsection{Industrial problem}
Conformal coatings are thin protective polymer films placed over electronic assemblies to mitigate moisture, contaminants and related electrical-reliability risks. Their optical transparency makes ordinary visible-light inspection difficult. Where the formulation contains an appropriate fluorescent tracer, illumination in the ultraviolet may reveal coating coverage. However, \emph{visibility of fluorescence is not a direct demonstration of acceptable coating thickness, adhesion, curing state or long-term protection}. IPC-HDBK-830A addresses coating selection, processing, coverage and inspection; IPC-A-610J addresses assembled-board acceptability, including coating coverage and thickness \citep{ipchdbk830a,ipca610j,ipccc830c}. The publicly accessible standard records were reviewed at the level of their tables of contents and summaries. They are not used here to invent numerical acceptance rules.

The target is a nominal 70~$\times$~70~mm populated PCB in repeatable production positioning. A camera viewing an approximately 80~$\times$~80~mm area is a tentative optical design assumption, providing a 5~mm margin around each nominal board edge. A fixed single-camera station may be materially simpler than scanning an entire panel with a moving XY platform. But this does not guarantee detection on component sidewalls, under tall components or behind connectors. Commercial systems explicitly add oblique cameras for those areas \citep{goepelcci,mirteccci}. Accordingly, the research concerns \emph{measurable visible regions under defined optical conditions}, not universal three-dimensional coating verification.

\subsection{Research aim and objectives}
The overarching aim is to assess whether a low-cost, fixed-geometry UV fluorescence station combined with reference-aware image analysis can perform reliable, auditable inspection of selected conformal-coating defects within a production-relevant cycle time. The measurable objectives are to: (i) characterize excitation, emission, camera response and registration error; (ii) define a board-specific coverage map from CAD, the coating plan and expert-labelled visible zones; (iii) compare rule-based and learned segmentation under identical conditions; (iv) identify and size defects with explicit keep-out/required-zone semantics; (v) quantify failure to detect actionable defects; and (vi) measure practical CPU-only end-to-end latency and cost.

A scientifically responsible claim requires a relevant comparator. Prior UV bubble inspection studies already combine classical candidate selection with CNN confirmation \citep{lee2021cnn,lee2021transfer}. Recent studies optimize UV intensity using statistical features and YOLOv8 tiling \citep{choi2024uv}, or train detector networks for dewetting and capillary defects \citep{patel2026coverage}. Therefore \emph{neither hybrid image processing nor CNN use for conformal coatings is automatically novel}. Novelty, if any, must arise from an isolated and measured property such as reference-constrained defect handling, illumination robustness, uncertainty-aware human review, resource efficiency, or validation across manufacturing lots.

\subsection{Research questions and falsifiable hypotheses}
\textbf{RQ1:} Under which controlled lighting and camera settings can UV fluorescence distinguish coating-covered from uncovered \emph{visible} PCB regions with repeatable intensity and acceptable registration error?\par
\textbf{RQ2:} Does a reference-aware hybrid improve \emph{board-level actionable-defect sensitivity at matched false-reject rate} relative to classical thresholding and standalone segmentation?\par
\textbf{RQ3:} Does multi-angle excitation or exposure bracketing reduce misses in shadow-prone regions enough to justify additional exposure/cycle-time cost?\par
\textbf{RQ4:} Can a CPU-class computer support the complete acquisition-to-result workflow within an actual specified production takt requirement, while meeting lighting stability and repeatability requirements?\par
\textbf{RQ5:} How do performance and calibration shift across physical boards, coating batches, operator handling, camera positioning and lighting aging?\par

\textbf{H1 (to test):} a geometric coating plan constrains false alarms from naturally dark components and improves protection of keep-out zones at fixed sensitivity. \textbf{H2 (to test):} flat-field normalization decreases performance variability across illumination geometries. \textbf{H3 (to test):} a lightweight segmentation network with reference-aware rules can be more resource-efficient than a substantially larger end-to-end model. The null hypotheses state that no meaningful improvement exists. Thresholds, comparator configurations and decision criteria must be frozen before evaluating the independent test set.

\subsection{Claims explicitly outside scope}
This manuscript does not assert full-coating thickness metrology by UV brightness; does not equate an image pixel with optical resolving power; does not infer that a 0.1~mm physical void can be detected reliably simply because it spans multiple pixels; and does not estimate a production false-acceptance rate from synthetic images. Destructive or OCT thickness measurements and electrical testing can independently validate selected hypotheses. OCT studies explicitly show how layer interfaces and thickness can be measured when fluorescence alone cannot provide such evidence \citep{shao2019oct,lf2023oct}.

\section{Literature search strategy and evidence appraisal}
\subsection{Search scope and transparency}
A preliminary structured source search was conducted on 23 October 2026 using public scholarly and official web records. Search clusters included \texttt{conformal coating PCB bubble CNN UV}, \texttt{coating Faster R-CNN dewetting}, \texttt{UV fluorescence optical filtering}, \texttt{IPC coating coverage standards}, \texttt{industrial CCI camera multi-angle CAD}, \texttt{PCB anomaly dataset}, \texttt{registration threshold semantic segmentation} and \texttt{OpenVINO CPU}. Publisher or institutional records were preferred for bibliographic identity, official dataset records for dataset access/licensing, standards-organization pages for standard versions, and original manufacturers for product claims. Search queries and limitations are stored in \texttt{sources/search\_log.csv}. Source records and review status are in \texttt{sources/source\_catalog.csv}.

This work is \emph{not yet a completed PRISMA systematic review}: a search-engine retrieval process was used to identify candidates, no complete database-specific result exports were collected, and not all full texts were acquired. PRISMA 2020 establishes useful expectations for transparent inclusion, screening and reporting \citep{page2021prisma}. The next academically rigorous stage is to repeat explicit query strings in IEEE Xplore, Scopus/Web of Science or the university's licensed databases, export results, deduplicate them, record decisions and then perform full-text appraisal. This document does not report a fictitious search flow or invented exclusion counts.

\subsection{Relevance hierarchy}
Evidence was classified into six levels: \textbf{direct coating-AOI research} (actual PCB conformal coating under relevant optics); \textbf{coating metrology} (OCT and other independent thickness approaches); \textbf{commercial coating inspection} (product capability claims); \textbf{transferable PCB inspection} (copper, assembly or solder defects); \textbf{general segmentation/anomaly algorithms and datasets}; and \textbf{standards, optics and system safety}. A paper testing ordinary PCB copper-trace defects is not experimental confirmation of UV-fluorescent coating defect detection. Similarly, the performance of Mask2Former on ADE20K or EfficientAD on MVTec cannot be used as the expected performance of a 70~mm coated PCB \citep{cheng2022mask2former,batzner2024efficientad,bergmann2019mvtec}.

\subsection{Evidence quality and extraction protocol}
The systematic full-text review should extract: imaging target; coating and tracer chemistry; wavelength and optical filter; FOV, image resolution and lighting geometry; physical specimen and production-lot counts; defect class definition and independent confirmation; model and hyperparameters; augmentation; leakage-free splits; measured defect-level and board-level metrics; false accept/reject counts; hardware/timing boundary; and external generalization. If any field is absent it must remain \emph{not reported}; do not infer a sample size from the number of augmented image patches. Distinguish:\par
\begin{enumerate}[leftmargin=*]
\item \textit{Full accessible publisher HTML or full paper examined}: explicit evidence may be extracted from results and methodology.
\item \textit{Publisher abstract / bibliographic record examined}: only its publicly exposed assertions may be recorded.
\item \textit{Vendor specification}: capability claims require independent confirmation before they become scientific evidence.
\item \textit{Standard preview or table of contents}: used to establish relevance and revision, not to reproduce protected acceptance tables.
\end{enumerate}

\subsection{Limitations of this preliminary corpus}
The source catalog contains a sizeable cross-disciplinary starter collection, but the relative number of dedicated coating papers is far smaller than the number of general AI papers. Some directly relevant studies are published in Korean, and detailed methods may require translation and rights-respecting full-text access \citep{lee2021cnn,choi2024uv}. Publication timing also matters: a conference presented in 2025 but published as a 2026 Springer chapter is cited using the confirmed 2026 publication record \citep{patel2026coverage}. The absence of a verified public dataset in the searched sources must be expressed as \emph{not located in the present search}, not as proof that no dataset exists anywhere.

\subsection{Operational research matrix}
The included machine-readable evidence matrix is intentionally conservative. Its \texttt{source\_verification} column should be upgraded from bibliographic/abstract verification only when methods, figures and supplementary datasets have actually been read. Each drafted numerical statement should include a \texttt{reported\_metric\_definition} field: detection recall, image classification accuracy, pixel IoU and manufacturing acceptance are different quantities. Before publication, a second reviewer should spot-check all references and all extracted performance values, ideally blinded to method ranking.

\section{Conformal coatings, defects and fluorescence optics}
\subsection{Protection mechanisms and material variability}
Conformal coating denotes a thin film that follows an electronic assembly's geometry. IPC-HDBK-830A documents acrylic, epoxy, silicone, polyurethane and parylene systems, along with application, curing, compatibility, coverage and thickness considerations \citep{ipchdbk830a}. These chemistry families differ in physical properties and processing requirements; even within a nominal material family, formulations, fluorescent tracer concentration and fluorescence spectra may vary. The production coating data sheet and qualification results are therefore indispensable experimental inputs. A coating that does not provide sufficient fluorescence under 365~nm irradiation invalidates the chosen imaging modality and must trigger a design review rather than post hoc enhancement of unreliable images.

A workable inspection taxonomy must separate physical mechanism from image appearance. \textit{Missing coating} denotes uncovered required material. \textit{Overspray or keep-out violation} means a forbidden location shows coating evidence. \textit{Bubble} means trapped gas in or under a film; \textit{blister} may denote raised separation of coating from substrate. \textit{Dewetting} creates locally receded coating and boundaries; \textit{capillary pooling} accumulates film near leads and geometrical gaps. \textit{Cracks, wrinkles, scratches and contamination} require specific visible criteria and possibly alternate lighting. These classes should be treated as candidates for annotation, not guaranteed detectable classes. IPC-A-610J includes conformal-coating coverage and thickness categories, and IPC's 2024 release specifically noted revisions clarifying coating voiding/bubbles \citep{ipca610j,ipc2024release}. The applicable class, customer drawing, and actual acceptance language must be obtained by authorized access to the complete requirements.

\subsection{Excitation, emission and optical contrast}
A fluorescent indicator absorbs excitation radiation and emits longer-wavelength light. In a simple first-order model, the measured camera signal at board position $\mathbf{x}$ can be written
\begin{equation}
I(\mathbf{x}) \approx g\, t\, q(\mathbf{x})\,E_{365}(\mathbf{x})\,\eta(\mathbf{x}) + B(\mathbf{x}) + \epsilon(\mathbf{x}),
\label{eq:fluorescence}
\end{equation}
where $g$ is camera/system gain, $t$ the exposure, $q$ the combined camera/lens/filter response, $E_{365}$ local excitation irradiance, $\eta$ an effective emission response of tracer and coating geometry, $B$ residual background/reflections and $\epsilon$ noise. Equation~\ref{eq:fluorescence} is an engineering abstraction, \emph{not} a calibrated film-thickness law. In particular, $\eta$ may depend nonlinearly on film thickness, concentration, photobleaching, shielding, topography and chemical condition. Brightness differences can therefore indicate suspicious coating variation but cannot, without independent characterization, prove an absolute thickness.

The correct filter is determined by \emph{what the camera should measure}. The UV illumination excites the fluorescent coating, while the camera should ordinarily transmit the visible emission and suppress reflected 365~nm radiation. Crowther systematically examined excitation-light blocking and undesirable filter/lens fluorescence in UV photography \citep{crowther2022filters}. Fluorescence microscopy technical literature similarly distinguishes excitation and emission filtering \citep{zeissfluorescence}. A 365~nm passband filter mounted on the camera would pass excitation and could reject much of the fluorescence signal; it should not be selected automatically simply because the LEDs are 365~nm. The protocol must first measure the coating emission spectrum (or empirically compare candidate long-pass/band-pass emission filters) and image a bare PCB under identical illumination.

\subsection{Field of view, physical sampling and resolution}
With an 80~mm horizontal FOV and a 4000-pixel image width, image sampling is
\begin{equation}
s_x=\frac{80\,\mathrm{mm}}{4000}=0.020\,\mathrm{mm/pixel}=20\,\mu\mathrm{m/pixel}.
\end{equation}
A 0.20~mm geometric feature covers about 10 pixels horizontally in the ideal projected image. But actual resolvability is constrained by modulation transfer function, lens focus/aperture, spatial contrast, diffraction, motion blur, photon statistics, demosaicing and the segmentation threshold. Furthermore, a 0.20~mm \emph{uncoated region} and a 0.20~mm \emph{bubble embedded under a fluorescent layer} may have different visibility. The smallest reliably detectable defect must be measured across known-size references and background conditions. The board's 70~mm dimension does not imply that a 12~MP camera is necessarily required; 5, 8 and 12~MP alternatives should be compared against measured critical defect sizes and optical contrast.

\begin{table}[htbp]\centering\small
\caption{Approximate horizontal spatial sampling for an 80 mm field of view; not a claim of defect-detection capability.}
\begin{tabular}{lrrr}\toprule
Image width & Sampling ($\mu$m/px) & 0.10 mm feature (px) & 0.20 mm feature (px) \\\midrule
1600 px & 50.0 & 2.0 & 4.0\\
2448 px & 32.7 & 3.1 & 6.1\\
3264 px & 24.5 & 4.1 & 8.2\\
4000 px & 20.0 & 5.0 & 10.0\\
\bottomrule\end{tabular}
\label{tab:sampling}
\end{table}

\subsection{Lighting geometry and experimental controls}
Component heights cause bright specular regions, self-shadowing and occlusion. A top-only flood source may illuminate unobstructed PCB areas consistently but miss coating on vertical component faces. Side UV sources can recover some information at the expense of exposure count or simultaneous irradiance complexity; commercial designs use both top and oblique views \citep{goepelcci,mirteccci}. The direct Choi et al. study is highly pertinent: it uses brightness/hue features, Mahalanobis distance and YOLOv8 tiling to evaluate appropriate UV intensity for bubble detection \citep{choi2024uv}. That prior work removes any basis for claiming UV exposure optimization in general as unprecedented.

A reproducible optical study should acquire a dark frame and an image of a stable fluorescence reference at each lighting condition. Exposure/gain and lighting must be logged. Temperature and UV source warm-up should be controlled; repeated references can quantify drift. Four-source geometry should initially be treated as an experimental independent variable (left/right/top/bottom or quadrants) rather than as a fixed recommendation. A simple single UV exposure may outperform a multi-exposure fusion system in throughput, if its shadows are tolerable for the target defect classes.

\subsection{Thickness versus fluorescence and independent validation}
The ability to image coverage is not equivalent to measuring coating thickness. Shao et al. measured the interfaces of conformal coatings non-destructively using ultra-high-resolution optical coherence tomography \citep{shao2019oct}. A later low-cost line-field Fourier-domain OCT study demonstrated a three-dimensional measurement on a coated PCB and reported, for one illustrated sample, approximately 27.2~$\pm$~7~$\mu$m on resistors versus 82.9~$\pm$~13~$\mu$m on the solder mask \citep{lf2023oct}. Those are results of OCT measurements on specific samples, not expected thickness values for this project. OCT is a useful \emph{ground-truth or secondary metrology instrument}, particularly in zones where fluorescence cannot distinguish merely thin from absent coating. If access to OCT is unavailable, carefully calibrated microscopy or an appropriate coating-thickness gauge may serve as independent evidence on test coupons.

\subsection{Safety constraints for optical development}
365~nm UVA can present photobiological hazards to eyes and skin. IEC~62471 provides classification and measurement concepts for optical-radiation hazards from lamps including LEDs \citep{iec62471}; ISO~12100 sets out machine risk assessment and risk-reduction principles \citep{iso12100}. Neither a nominal LED optical wattage nor a wavelength alone determines safe exposure. The prototype shall use an opaque UV enclosure, interlocked access, power-off-on-fault behavior and documented verification of leakage and service states. This report provides research requirements, not certification of machine safety.

\section{Direct prior art and industrial inspection systems}
\subsection{Direct academic research on coated PCBs}
The most relevant academic evidence is much narrower than the general literature on bare-PCB surface defects. The two 2021 Korean Journal of Broadcast Engineering papers by Lee et al. directly address conformal-coated PCB bubble inspection \citep{lee2021cnn,lee2021transfer}. In the earlier approach, UV images were processed to find dark-centre and bright-boundary candidates, after which CNN-based verification rejected non-bubble lookalikes. The authors' abstract reports precision 0.805, recall 0.763 and F1 0.767. As those rounded values are not algebraically consistent with the conventional harmonic-mean F1 calculated from precision and recall, the reported metric definitions and aggregation method require inspection of the complete paper before they are used as a benchmark. This discrepancy is precisely why the original evidence matters. In the later paper, VGG-style CNN transfer learning is reported to improve bubble detection to F1 0.9044. Those results cannot be directly compared without matching data, test splits and the exact positive class.

Zouhri and colleagues presented a Faster R-CNN assessment of coating bubbles at ICM 2021 and published a related 2024 journal discussion of machine-learning and edge inspection \citep{zouhri2021faster,zouhri2024ai}. Their presence further establishes the pre-existing literature on deep-learning-based coating detection; however the available bibliographic record is insufficient to evaluate the size of their dataset or reproduce their results. They remain marked for full-text follow-up.

In 2024, Choi et al. explicitly addressed UV irradiance selection for small conformal-coating bubbles, using image-level colour/brightness descriptors, feature-space distances, YOLOv8 and image tiling \citep{choi2024uv}. This is essential to the lighting and ablation design: it means our proposed illumination normalization must be positioned relative to existing UV intensity suitability approaches. Wang et al. reported a YOLOv8-Nano variant with attention and feature fusion for PCB coating bubble detection, with a 12\% relative performance improvement for smaller bubbles in the institutional abstract \citep{wang2025bubble}. That figure is \emph{source-reported and method-specific}, and cannot be treated as an absolute accuracy gain or as performance on the present boards.

More recently, Patel et al. presented a customised Faster R-CNN with ResNet-152 and enhanced proposal mechanisms for capillary and dewetting patterns on conformal-coated PCBAs \citep{patel2026coverage}. The publisher/institutional abstract reports recall of 98.53\% and 95.49\% for two defect categories, respectively. This is relevant because it moves past bubbles toward process-related coverage flaws, but recall alone says nothing about precision, false rejects, specimen independence or line deployment. A 2026 hybrid PCB inspection chapter also reports a combination of registration, component-fault learning and HSV analysis of \emph{excess conformal coating} \citep{pidurkar2026hybrid}. Its presence means the novelty of a generic classical-plus-AI pipeline must be framed very narrowly.

\subsection{Relevant adjacent work}
General PCB AOI surveys discuss sensor/illumination choices, geometric comparison, segmentation, classification and manufacturing integration \citep{ebayyeh2020aoi,zhou2023pcb,ieee2023pcb}. A 2025 PCB fault system aligns boards using RANSAC and segments defects using U-Net under several illumination conditions \citep{qub2025registration}. This is useful for designing the \emph{geometric and lighting baselines}, but its conventional PCB defects cannot validate coating-coverage decisions. Similarly, the old 1988 Hara et al. system exploited ultraviolet-excited \emph{PCB base-material fluorescence} to distinguish copper patterns \citep{hara1988fluorescent}; it is historical evidence for fluorescence-based imaging in PCB inspection, not evidence for modern coating evaluation.

Outside electronics, Han et al. used autoencoders, self-labelling, and convolutional models to identify coating boundaries on fuel-injection nozzles \citep{han2025coating}. Their fully accessible paper reports over 95\% interface localization within $\pm 6$ pixels and 7.18 images per second, along with important limitations on generalization to other sources. The concept of reducing manual annotation may transfer, but its sample geometry, surface appearance and lighting differ markedly from densely populated PCBs. Broader coating-defect surveys likewise provide method categories rather than transferable PCB inspection performance \citep{zi2026coating}.

\subsection{Commercial conformal-coating AOI systems}
\begin{table}[htbp]\centering\footnotesize
\caption{Published manufacturer descriptions. Numeric specifications are vendor claims, not independent head-to-head benchmarks.}
\begin{tabularx}{\textwidth}{p{23mm}X X}\toprule
Supplier / product & Relevant published capabilities & Engineering implication \\\midrule
Nordson FX-940UV \citep{nordsoncci} & 5 MP colour image, UV illumination, taught known-good coating/noncoating areas, optional oblique views; manufacturer describes inspection in a few seconds. & Golden reference plus rules are established practice; 5 MP can be commercially relevant for some configurations.\\
G\"opel Multi Line CCI \citep{goepelcci} & Top and side UV, CAD/coating-plan regions, claimed 9/17 $\mu$m sampling, optional multiple oblique cameras and separate thickness measurement. & Region semantics and geometric coverage plans are established; complex 3D areas demand more than one view.\\
MIRTEC GENESYS-CC \citep{mirteccci} & 50 MP camera, claimed 4 $\mu$m optics and 30 $\mu$m microbubble sensitivity; AI-based bubble assessment, additional UV and side camera geometry. & Defect size and inspection risk determine resolution; published claims need independent validation.\\
Viscom CCI \citep{viscomcci} & UV-triggered fluorescence for coating presence/evenness; distinguishes optical inspection from spot thickness measurement. & Treat fluorescence intensity as a process signal, not calibrated thickness without metrology.\\
\bottomrule\end{tabularx}
\label{tab:commercial}
\end{table}

The existence of these systems has two consequences. First, CAD-driven forbidden/required regions, UV illumination and multiple camera views cannot individually be claimed as new. Second, commercial product brochures rarely provide a publicly reproducible dataset with independent board counts and false-acceptance operating points. Academic work can add value by publishing \emph{transparent evaluation and cost-performance trade-offs} rather than simply reproducing a proprietary feature list. A credible comparison should not equate our €-level prototype estimate to a large commercial machine including conveyors, warranty, service, traceability and optional three-dimensional metrology.

\subsection{Synthesis: what has not been established by available evidence}
Direct papers demonstrate the feasibility of learning bubble and selected coverage-defect patterns, but they do not collectively establish an openly benchmarked low-cost, single-view, CPU-only system meeting a specified factory false-acceptance rate across independent coating lots. The reviewed commercial pages demonstrate use of reference plans, but do not independently establish how much they improve each specific defect class. Thus an honest research opportunity is to isolate the \emph{incremental value of reference awareness and illumination correction} under controlled tests. It remains possible that a calibrated classical thresholding algorithm will outperform a neural network for simple missing-coating or keep-out checks, while neural processing is needed only for complicated bubbles and texture defects.

\section{Algorithms, public datasets and methodological trade-offs}
\subsection{Classical vision is a necessary baseline}
For a single known PCB design under controlled UV illumination, classical vision offers a strong, interpretable starting point. Correct dark frames, gain and spatial irradiation; align to fiducials or stable geometric features; threshold a normalized coating-likelihood channel; apply small, physically justified morphological operations; and compare the resulting film mask with required and forbidden polygons. Otsu's histogram thresholding is a principled unsupervised baseline \citep{otsu1979}. Geometric alignment can use fiducials and a planar projective transformation with robust RANSAC fitting \citep{fischler1981ransac}; camera intrinsics and lens distortion can be estimated with planar calibration \citep{zhang2000calibration}. Image registration reviews distinguish feature-based from intensity-based matching and the risks from resampling/occlusion \citep{zitova2003registration}. On boards with rich component textures the reference image may be stable; when strong UV appearance changes make feature matching unreliable, acquire an auxiliary white-light registration image or use mechanical fiducials.

The main classical failure modes are nonuniform excitation, glossy reflectance, PCB fluorescence unrelated to the coating, component shadows, natural variation in tracer brightness and threshold sensitivity. Morphology can artificially erase a small genuine defect; the minimum structuring-element scale must be expressed in $\mu$m using measured image scale, and evaluated with defect-size-dependent recall. Rules should be engineered from independent validation data, not manually tuned on the final test boards.

\subsection{CNN and transformer segmentation}
U-Net remains a reasonable low-data segmentation comparator because its skip connections support preservation of spatial detail \citep{ronneberger2015unet}. DeepLabv3+ uses atrous processing and a boundary-refining decoder, offering a more complex reference model \citep{chen2018deeplab}. Lightweight backbones such as MobileNetV3 may reduce CPU cost \citep{howard2019mobile}. These methods were designed and first evaluated on other domains; UV coating segmentation performance must be learned empirically.

Transformers provide a further family rather than an automatic upgrade. Swin Transformer aggregates hierarchical contextual representations \citep{liu2021swin}; SegFormer pairs multi-scale features with a small MLP decoder \citep{xie2021segformer}; Mask2Former supports multiple mask segmentation tasks \citep{cheng2022mask2former}; MobileViT investigates low-resource global representation \citep{mehta2022mobilevit}. The real evaluation criterion here is not a published score on ADE20K, but sensitivity to relevant physical flaws, annotation requirements, false-reject behaviour, CPU latency and reproducibility. A practical progression is (1) a classical reference-aware baseline; (2) small U-Net; (3) SegFormer-B0 or MobileNetV3-based segmentation if justified; and only then higher-capacity networks if systematic errors warrant them.

For very small bubbles, downsampling a full 4000-pixel image to 512 or 640 pixels wide changes effective sampling from 20 to 125--156 $\mu$m/px. Consequently, a nominal 0.2~mm defect that spans ten original pixels may span only about 1.3--1.6 pixels after aggressive resize. The model evaluation must explicitly preserve resolution by using a high-resolution input, carefully selected tiles, a ROI hierarchy or a dedicated small-defect head. Tile counts and overlapping-region merging are part of runtime measurement. The Choi et al. paper is direct evidence that tiling has been investigated for tiny coating bubbles \citep{choi2024uv}; tiling as such is not novel.

\subsection{Anomaly detection and scarce defects}
Industrial inspection often has many acceptable specimens and few verified examples of rare defects. PatchCore represents acceptable image patches with a feature-memory bank \citep{roth2022patchcore}; DRAEM trains with synthetically generated anomalies \citep{zavrtanik2021draem}; EfficientAD explores a low-latency student--teacher arrangement and global context \citep{batzner2024efficientad}. The large accuracy and millisecond figures from those papers are benchmark- and hardware-specific. In coated PCB production, \emph{normal variation} includes acceptable fluorescence differences, changing component reflectance and permissible masking areas, so a normal-only model can flag harmless variations while missing subtle functionally meaningful deviations. If used, anomaly scores should be integrated with reference-region rules and evaluated at a selected board-level false-reject operating point.

The uncertainty-aware PCB method $\mathrm{U}^{2}\mathrm{D}^{2}$PCB illustrates reconstruction and discriminative anomaly logic on conventional PCB data \citep{u2d2pcb2024}. It is a useful architectural candidate for low-defect-frequency problems but no substitute for an actual UV coating dataset. Likewise, CM-UNetv2 targets small ordinary PCB features and boundaries, offering a hypothesis about feature preservation rather than evidence of coating behavior \citep{guo2025cmunet}.

\subsection{Dataset inventory and actual relevance}
\begin{table}[htbp]\centering\footnotesize
\caption{Verified public benchmark identities. None of these is a demonstrated public pixel-labelled UV PCB conformal-coating benchmark.}
\begin{tabularx}{\textwidth}{p{25mm}p{25mm}X p{27mm}}\toprule
Dataset & Scale / access & What it contains & Use in present study \\\midrule
DeepPCB \citep{tang2019deeppcb} & 1,500 template/test pairs; source released & Conventional PCB pattern defects, six categories & Registration/detection pipeline only\\
DsPCBSD+ \citep{lv2024dspcb,dspcbfigshare} & 10,259 images, 20,276 boxes; Figshare CC BY & Copper residue/deficiency, scratches and foreign objects; local crops & Annotation tooling, small-object benchmark\\
FPIC \citep{jessurun2022fpic} & Research collection described in preprint & Semantic optical PCB assurance images & General semantic PCB task exploration\\
MVTec AD \citep{bergmann2019mvtec,mvtecdataset} & 5,354 images; noncommercial license & 15 object/texture categories, anomalies & Anomaly-code sanity check only\\
MVTec LOCO AD \citep{bergmann2022loco} & 3,644 images, noncommercial & Structural and logical anomalies & Test reference-rule ideas conceptually\\
VisA \citep{visa2022dataset} & 10,821 images, CC BY & Industrial anomalies including PCB-like categories & Transfer learning or anomaly comparison, not coating ground truth\\
\bottomrule\end{tabularx}
\label{tab:datasets}
\end{table}

DsPCBSD+ provides explicitly downloadable images and object annotations under an official Figshare record \citep{dspcbfigshare}. This makes it useful for validating image loaders, annotation conversion and neural training code, but those image labels refer to bare PCB surface/copper defects. DeepPCB has reference/test image pairs and defect boxes, which could help test registration software before any UV data are captured \citep{tang2019deeppcb}. MVTec AD and LOCO AD carry non-commercial dataset licences; they must not silently migrate into a factory training or redistribution workflow \citep{mvtecdataset,bergmann2022loco}. VisA lists a more permissive CC BY dataset licence but is still not a substitute for collecting UV coating images \citep{visa2022dataset}.

\subsection{Dataset acquisition is a central experimental contribution}
A purpose-built UV coating dataset is therefore a \emph{necessary engineering deliverable} and possibly an independent scholarly contribution. It should include nominal 70~mm boards across physical board IDs, coating lots, component configurations and controlled acquisition conditions; carefully separated genuine defects and induced test coupons; expert-adjudicated defect masks or instances; and reference areas labelled as \emph{required, forbidden, permitted/optional, occluded and unknown}. In any published dataset, source images and annotations need a release decision independent of internal production images and CAD. One hundred repeat exposures from ten physical PCBs do not constitute one hundred independent defect specimens for statistical inference. Photometric augmentation cannot create genuinely new physical bubble morphologies or validate transfer to a novel coating batch.

\subsection{Prior-work synthesis matrix}
The detailed catalog in the accompanying CSV records 58 source identities and their verification status. The literature comparison in Table~\ref{tab:prior} summarizes why the architecture should be judged incrementally rather than advertised as a new generic fusion of known methods.
\begin{table}[htbp]\centering\scriptsize
\caption{Directness of the principal academic comparisons.}\label{tab:prior}
\begin{tabularx}{\textwidth}{p{28mm}X p{25mm}}\toprule
Study & Primary relevant contribution & Evidential constraint\\\midrule
Lee et al. 2021a/b \citep{lee2021cnn,lee2021transfer} & Coating bubble UV candidate extraction plus CNN / transfer learning & Defect type primarily bubbles; no demonstrated production transfer to our line\\
Choi et al. 2024 \citep{choi2024uv} & UV level suitability, multitype bubble location and tiled YOLOv8 & Does not settle other coating defects or CPU runtime\\
Wang et al. 2025 \citep{wang2025bubble} & Enhanced YOLOv8-Nano for small PCB coating bubbles & Relative improvement reported; detailed sample counts unverified\\
Patel et al. 2026 \citep{patel2026coverage} & Faster R-CNN for dewetting/capillary coating defects & Class recalls without general production false-acceptance operating point\\
Han et al. 2025 \citep{han2025coating} & Self-labeling coating-interface localization & Fuel-injection nozzles, not UV PCB surfaces\\
Shao et al. 2019 \citep{shao2019oct} & Independent non-destructive coating thickness metrology & Hardware and observable different from fluorescence coverage\\
\bottomrule\end{tabularx}
\end{table}

\section{Proposed reference-aware inspection architecture}
\subsection{Design rationale and data flow}
The proposed architecture intentionally does not require inventing a new backbone before observing failure cases. A constrained board-specific coating plan defines \emph{where} coating is required or forbidden; the optical model measures \emph{whether} sufficient fluorescence evidence is observed; and an optional learned branch estimates ambiguous coating structure or defect patterns. This separation provides traceable explanations for inspection decisions and facilitates change control when a PCB revision modifies the coverage plan. Figure~\ref{fig:architecture} presents the proposed data flow.

\begin{figure}[htbp]\centering
\begin{tikzpicture}[>=Stealth,node distance=0.70cm and 0.55cm,
 block/.style={draw,rounded corners,align=center,minimum height=0.85cm,text width=3.7cm,font=\footnotesize},
 wide/.style={draw,rounded corners,align=center,minimum height=.85cm,text width=6.7cm,font=\footnotesize}]
\node[wide] (acq) {Controlled 365 nm UV exposure(s) + optional white-light registration image};
\node[wide,below=of acq] (qc) {Acquisition QC: metadata, dark frame, exposure, saturation, fluorescent reference};
\node[wide,below=of qc] (geo) {Undistort + align board to coating-plan coordinates; registration quality gate};
\node[block,below left=.85cm and 0.1cm of geo] (model) {Classical / learned coating probability and possible defect confidence};
\node[block,below right=.85cm and 0.1cm of geo] (ref) {Approved reference map: required, forbidden, optional, occluded};
\node[wide,below=2.0cm of geo] (logic) {Area, boundary, connected regions, keep-out and critical-zone decision rules};
\node[wide,below=of logic] (risk) {Pass / fail / manual review; reason, overlay, confidence and archived images};
\draw[->] (acq)--(qc); \draw[->](qc)--(geo); \draw[->](geo)--(model); \draw[->](geo)--(ref);
\draw[->](model)--(logic);\draw[->](ref)--(logic);\draw[->](logic)--(risk);
\end{tikzpicture}
\caption{Hypothesized inspectable and auditable architecture. This diagram represents a proposed pipeline, not experimentally validated improvement.}\label{fig:architecture}
\end{figure}

\subsection{Radiometric correction and quality gate}
For an illuminated raw image $I$ with dark frame $D$ and stable flat fluorescent reference $F$ imaged under the same conditions, one possible corrected signal is
\begin{equation}
I_c(\mathbf{x}) = \frac{I(\mathbf{x})-D(\mathbf{x})}{\max(F(\mathbf{x})-D(\mathbf{x}),\epsilon)}\,\overline{(F-D)}.
\label{eq:flatfield}
\end{equation}
Here $\epsilon$ prevents division by near-zero values and $\overline{(F-D)}$ controls overall scaling. Correction is accepted only where the reference is stable, correctly registered and not clipped. It \emph{does not} unshadow occluded surfaces or compensate for all changes in the fluorescent tracer's chemistry. Quality gates flag saturation, insufficient brightness, excessive dark-frame residual, missing board, out-of-focus image and out-of-distribution lighting conditions. A failed acquisition gate must produce \emph{review/reacquire}, not an automatic PASS.

\subsection{Registration and reference coordinates}
Let $\mathbf{p}$ denote an image pixel and $\mathbf{q}$ its location in the coating-plan coordinate frame. A planar projective transform $\mathbf{q}\sim H\mathbf{p}$ can be estimated from fiducials or stable features using RANSAC \citep{fischler1981ransac,zhang2000calibration}. Where board geometry is rigid and viewing is normal, a similarity or affine transformation may be adequate and less prone to overfitting. Registration uncertainty should be measured on withheld fiducials or physical landmarks. A failed registration quality gate must not silently propagate a displaced keep-out mask to product disposition. Precision can be reported as median, 95th percentile and maximum pixel reprojection errors, converted to millimetres by a calibrated transform.

Approved mask regions form a disjoint partition of the visible board: $R$ = must-coat, $K$ = must-not-coat, $A$ = optional/allowed, and $U$ = not reliably inspectable/occluded. Every pixel belongs to exactly one region after resolving drawing overlaps. This explicit \emph{unknown/occluded} class prevents classification of hidden material as acceptable simply because the image contains no evidence. The plan revision is a versioned input, ideally derived from approved manufacturing data rather than an arbitrarily labelled ``perfect'' reference photo.

\subsection{Segmentation output and interpretable indicators}
Define $P(\mathbf{x})\in[0,1]$ as estimated probability or standardized evidence of coating, and $M_\tau(\mathbf{x})=\mathbb{1}[P(\mathbf{x})\geq\tau]$ as a binary mask chosen on validation data. Required-zone coverage and forbidden-zone contamination can be expressed as
\begin{equation}
C_R=\frac{|M_\tau\cap R|}{|R|},\qquad
C_K=\frac{|M_\tau\cap K|}{|K|}.
\end{equation}
These are \emph{image-derived ratios}, not coating-thickness measurements. Overall coverage may conceal a critical local hole; therefore identify connected missing regions $H_i\subset R\setminus M_\tau$, measure each area $A_i = n_i s_xs_y$, maximum inscribed radius, contour length and proximity to critical features. For forbidden areas, analogous connected coating evidence is measured. Add image confidence and consistency across illumination exposures. Rules can assign separate criticality to connector pins, pads, RF areas or other zones, but only when approved by the manufacturing engineer.

A nominal defect area threshold $A_{\min}$ should reflect actual product acceptance criteria and the optical measurement error. The decision logic may be written
\begin{equation}
D=\begin{cases}
\mathrm{REVIEW}, & Q=0 \text{ or uncertainty $>u_{\max}$},\\
\mathrm{FAIL}, & \exists\text{ validated actionable region violating $R$ or $K$ rules},\\
\mathrm{PASS}, & \text{otherwise within defined visible inspection scope},
\end{cases}
\end{equation}
where $Q$ summarizes acquisition, registration and model validity. A PASS must be explicitly scoped to inspected visible surfaces and validated defect categories; it must not imply that hidden component sides have been inspected.

\subsection{Model selection and ablation logic}
The segmentation branch should be replaceable: a calibrated HSV/intensity threshold, lightweight U-Net, MobileNetV3 decoder, SegFormer-B0 or a learned anomaly score. The reference-comparison branch must use identical region rules and optical inputs where possible across experiments. Ablations should remove individual modules one at a time: (A) flat-field correction; (B) geometric registration; (C) the learned mask; (D) the reference regions; (E) multi-angle illumination; and (F) uncertainty handling. If the hybrid's improvement disappears after the reference rules are applied to the classical baseline, the claimed innovation is better described as a useful industrial system design than as an ML architectural breakthrough.

\subsection{Implementation interfaces}
Recommended standalone software modules are \texttt{acquisition.py}, \texttt{calibration.py}, \texttt{registration.py}, \texttt{segmentation.py}, \texttt{reference\_map.py}, \texttt{defect\_metrics.py}, \texttt{decision.py}, \texttt{traceability.py} and \texttt{benchmark.py}. Each module should expose typed inputs, version, provenance and error states. Store the model weights, preprocessing parameters, mask revision, threshold revision, image hashes, human review and machine decision. For factory deployment, a deterministic result protocol and watchdog are as important as the selected model. ONNX Runtime's execution-provider model and current OpenVINO CPU support provide plausible portability options, but not guaranteed timing on an unspecified Intel PC \citep{onnxproviders,openvino2026}.

\section{Reproducible experimental methodology}
\subsection{Study design and staged feasibility gates}
The proposed research is a prospective, controlled comparative study. Gate~0 establishes whether the actual coating emits measurable light at the proposed excitation wavelength, and whether the UV enclosure is safe. Gate~1 establishes optical repeatability and ground-truth reproducibility. Gate~2 develops three algorithmic approaches with separate training and validation boards. Gate~3 freezes all thresholds and models, then measures performance on completely held-out boards and ideally held-out production lots. Gate~4 evaluates production robustness, human override and timing over a sufficiently long representative sequence. Every stage has an explicit failure condition: if critical defects are not optically observable, further neural training alone is not a satisfactory resolution.

\subsection{Specimen and lot sampling}
Each physical PCB must have a unique immutable \texttt{specimen\_id}; each production/process batch has \texttt{lot\_id}. Every image receives an \texttt{image\_id} linking it to coating material, preparation, light arrangement, camera parameters, board orientation, cure status, capture time and annotation revision. Acquire multiple images from the same physical board to test optical repeatability but keep all images from that board \emph{within one data partition}. Where enough lots exist, use grouped splits by lot (train lots, validation lots, independent test lots). If only one lot exists, report that generalization across lots has not yet been tested. The companion \texttt{dataset\_register.csv} provides required tracking columns.

An initial \emph{planning target} could be approximately 100--200 physical boards across at least three visually distinct collection sessions or lots; this is not a power calculation and should not be presented as automatically sufficient for a rare-defect detection claim. Exact physical board counts must follow expected defect prevalence, the minimum acceptable sensitivity, the feasible number of independently confirmed defect boards, and confidence-interval targets. If natural faults are rare, create deliberately defective test boards with controlled omission, overspray, bubble-generating process variation or masking failures, subject to manufacturing and safety approval. Artificial examples should supplement, not replace, naturally occurring failures in the final production validation. A synthetic pixel hole applied in software is a useful stress test but not a true coating defect specimen.

\subsection{Defect taxonomy and acceptance thresholds}
\begin{table}[htbp]\centering\footnotesize
\caption{Candidate annotation and decision ontology. Inspection feasibility and acceptance limits are to be established, not assumed.}\label{tab:ontology}
\begin{tabularx}{\textwidth}{p{30mm}p{25mm}X}\toprule
Class & Label representation & Independent confirmation / concern\\\midrule
Missing coating & Pixel polygon or region instance & UV evidence plus expert/alternate inspection; distinguish intentional mask-out\\
Excess coating / overspray & Pixel region crossing forbidden zone & Coating plan revision and expert confirmation\\
Bubble / void & Instance region, size and location & Microscopy or expert adjudication; fluorescent dark spots are not all bubbles\\
Dewetting & Boundary polygon, area and context & Process engineer confirmation; often resembles absence or capillary effect\\
Capillary pooling & Region plus component proximity & Fluorescence intensity alone may not quantify thickness\\
Crack / scratch & Instance or line segment & Alternate-angle white/UV imaging; distinguish surface artefacts\\
Occluded / not observable & Exclusion mask & Must be handled as unknown, not labelled sound coating\\
\bottomrule\end{tabularx}
\end{table}

A defect is \emph{actionable} only if a documented customer, internal drawing or authorized acceptance specification makes it rejectable under the applicable product class. Record minimum size, permitted locations, any exceptions for component leads and acceptable bubble conditions. The goal is traceable compliance with a validated inspection plan, not the maximum number of bright/dark pixels detected. The full IPC-A-610J and coating plan must be consulted by a suitably qualified person before production rules are finalized \citep{ipca610j,ipc2024release}.

\subsection{Optical acquisition protocol}
For each run, record LED electrical settings and ideally calibrated irradiance (or at least an inter-session stable radiometric monitor), exposure time, analog gain, white balance state, image bit depth, lens aperture/focus, working distance, board-to-camera distance and filter identification. Use lossless raw or minimally processed image data wherever practical; JPEG compression can distort tiny bubble contours. Capture a dark frame and fluorescence reference after equipment warm-up. Image a bare uncoated board or verified keep-out surface to estimate background auto-fluorescence. Confirm that the visible-emission band is dominant after optical filtering \citep{crowther2022filters}.

The nominal factorial optical study should vary lighting configuration (top-only, two opposing sides, four-source combined, sequential directional exposures), exposure (at least three levels avoiding saturation), and image scale (for example ~5, ~8, ~12 MP optical equivalents or calibrated downsampled conditions). Do not use a downsampled 12~MP image as complete substitute for testing an actual 5~MP sensor: lens pixel size, color filter array, noise and readout differ. Each configuration should include multiple independent boards in randomized capture order; sources of non-independence should be modelled or accounted for when estimating uncertainty. Full-factorial designs may be expensive; staged screening followed by a focused confirmatory design is acceptable if predeclared.

\subsection{Ground truth and observer agreement}
Ground truth must be independent of the algorithms being evaluated. Two trained inspectors should independently label a predefined subset, blinded to model output, then reconcile disagreement using qualified engineering or microscopy review. Store \emph{original labels}, adjudicated labels, defect class, criticality, size and ``uncertain'' states. Record inter-annotator agreement on actionable defect identification and segmentation where meaningful; IoU/Dice between masks and agreement on board disposition measure different things. Some ground-truth evidence may be obtained by selective re-inspection at higher magnification, microscopy or OCT. If two experts cannot reliably agree on an image, model failure against one arbitrary mask should not be misinterpreted as a purely algorithmic problem.

\subsection{Classical baseline A}
Develop a fully documented corrected-image threshold pipeline using Otsu/adaptive thresholds as alternatives, threshold channel selected using validation data, fixed morphology based on millimetre-scale rules, and connected-component grouping \citep{otsu1979}. Register the board geometry using fiducial-driven affine/projective fitting and fail closed on poor registration \citep{fischler1981ransac}. Evaluate both a \emph{simple threshold-only} system and a \emph{classical plus reference masks} system. Without this latter baseline, improved performance of the hybrid could merely come from adding the reference plan rather than introducing AI.

\subsection{Learned baseline B}
Start with a reproducible U-Net variant and, if needed, SegFormer-B0 or MobileNetV3-based segmentation \citep{ronneberger2015unet,xie2021segformer,howard2019mobile}. Record fixed random seeds, model input resolution, weight initialization, class balance, augmentation, loss definition, optimizer, learning rate and early stopping. Set architecture and input resolution before viewing held-out test data. Training data may be augmented for lighting shifts, rotations and focus/noise within realistic ranges; prohibit spatial warps that invalidate defect sizes or move keep-out masks independently of the image. Use the same physical-board split and comparable augmentation conditions for competing models.

\subsection{Hybrid candidate C and ablations}
Run exactly the same optical acquisitions through registration, a learned coating probability map and geometry-aware required/forbidden region analysis as proposed in Section~6. Predeclare rules for a thresholded missing-coating area, largest missing island, proximity to critical geometry and uncertain-region abstention. Experimental factors are the presence or absence of correction, registration, reference map, learned branch and alternate illumination. A useful factorial ablation table includes: classical only, classical+reference, CNN only, CNN+reference, CNN+reference+correction, and model with multiple exposures. Report where improved results arise from a particular stage. A learned feature extractor could be replaced with PatchCore or EfficientAD in a secondary comparison, but avoid an unmanageable unprioritized model zoo \citep{roth2022patchcore,batzner2024efficientad}.

\subsection{Defect localization and fair matching}
Set a physical detection threshold. For a predicted defect instance $P_i$ and independently annotated ground-truth instance $G_j$, match using intersection-over-union (IoU) threshold, centroid tolerance or any-domain-specific criterion chosen \emph{a priori}. Tiny line defects may need a tolerance in millimetres rather than IoU alone. Use one-to-one matching between predictions and ground truth to avoid double counting. Summarize counts by defect class, size, location, illumination and criticality. Count a missed actionable physical defect as a false negative even if the image classifier correctly identifies another defect elsewhere on the same board.

\subsection{Production timing and system verification}
The complete cycle includes trigger-to-exposure, optional multiple exposures, camera transfer, acquisition validation, geometric registration, any inference tiling, defect aggregation, logging/database output, PLC result handshake and safety/fault reporting. Measure median, 95th and 99th percentile times across sustained runs, along with cold start, model reload, CPU load, RAM, thermal throttling and recovery after camera disconnection. A nominal inference-only speed on a GPU-equipped laptop cannot be copied to a used Intel office PC. The workstation's actual takt requirement must be supplied by production engineering and the system benchmarked against it \citep{onnxproviders,openvino2026}.

\subsection{Reproducibility and leakage checks}
Before each experiment: (1) confirm dataset hash and physical-board split; (2) confirm all label revisions; (3) record software Git commit, Python/CV/ML library versions and hardware; (4) freeze train/validation/test membership; (5) save calibrated acquisition parameters; (6) archive complete confusion matrices and false-decision images. Training examples acquired from the same specimen as the test image are an especially likely source of inflated accuracy. Avoid using a test specimen as the reference template for its own detection or including its other exposures in training.

\section{Statistical analysis, reliability and production feasibility}
\subsection{Primary and secondary outcomes}
The proposed \emph{primary} outcome is board-level sensitivity for \emph{actionable} defects at a predeclared maximum false-reject rate on independent boards. For a binary decision, define defect-positive ground truth, true positives (TP), false negatives (FN), false positives (FP) and true negatives (TN) at the \emph{board level}. Then
\begin{equation}
\mathrm{Sensitivity}=\frac{TP}{TP+FN},\quad \mathrm{Specificity}=\frac{TN}{TN+FP},\quad
\mathrm{Precision}=\frac{TP}{TP+FP},\quad \mathrm{FPR}=\frac{FP}{TN+FP}.
\end{equation}
False acceptance of a physically defective board is the risk-relevant $FN$ event; false rejection of a conforming board is $FP$. These conventions must be attached to every report to avoid ambiguous ``false pass'' terminology. At defect level, report size-stratified recall with an explicit matching rule. At pixel level, report Dice and IoU but interpret them as supporting localization metrics, not the factory pass criterion. Report uncertainty alongside proportions and the exact numerator and denominator.

\subsection{Confidence intervals and sample-size planning}
Use a binomial Wilson or exact interval for board-level proportions when board outcomes are sufficiently independent. At low defect prevalence, an apparent ``100\% recall'' on only ten defective boards says little about industrial reliability. For an illustration, if zero false accepts occur among $n$ independent actionable-defect boards, a rough one-sided 95\% upper bound on miss probability is $3/n$ (the ``rule of three'' approximation). Therefore even zero misses among 100 defective boards supports only an upper bound near 3\%, \emph{not} a proven $<0.1\%$ miss rate. Correlation within coating lots lowers effective information, so stratify and report between-lot results rather than treating thousands of crops as independent trials. Derive final sample sizes from desired confidence limits and the process's defect distribution before testing.

\subsection{Comparing algorithms fairly}
All methods must evaluate the same independent boards under the same acquisition conditions and specified threshold policy. For binary paired board outcomes, a paired disagreement analysis (e.g. McNemar-type test where assumptions apply) is preferable to treating the models' performance estimates as unpaired. For localization IoU or cycle-time measurements, consider paired bootstrap resampling at the \emph{physical-board level}; if lot clustering is material, resample lots and boards appropriately. Primary tests and alpha level should be prespecified. Exploratory comparisons across dozens of architectures are explicitly labelled exploratory, with appropriate caution regarding multiplicity. Report effect sizes and confidence intervals even when formal significance is inconclusive. Save the list of individual failure cases, not just summary percentages.

\subsection{Repeatability and domain shift}
A repeated-measures study should test the same board at multiple positions, orientations and times. Separate within-session repeatability from between-day, lighting-source aging and coating-batch reproducibility. Train/validation/test splits by \emph{physical board} remove identical-image leakage, but they do not independently test future batches; a deliberately held-out lot is stronger. Apply structured perturbations to registration, illumination and exposure and measure performance degradation with a documented acceptable degradation threshold. If a new PCB revision changes component heights or mask geometry, perform change-impact analysis rather than assuming old models generalize. Record ``review'' decisions separately from false rejects, since a system that abstains excessively may not meet line takt.

\subsection{Failure-mode and risk analysis}
\begin{table}[htbp]\centering\footnotesize
\caption{Potential failure modes and the proposed monitoring response.}\label{tab:risk}
\begin{tabularx}{\textwidth}{p{28mm}X p{35mm}}\toprule
Failure & Consequence & Control / observable \\\midrule
LED power or spectral drift & Weak fluorescence appears as missing coating & Stable reference target, drift alarm, locked settings\\
Wrong/uncalibrated emission filter & Reflected UV dominates signal & Spectral check, reference subtraction, qualification records\\
Board misregistration & False keep-out failures or masked true flaws & Fiducial residual threshold, manual review state\\
Specular component/occlusion & Missing sidewall or dark inaccessible areas & Exclusion/occlusion map, optional side views\\
Coating batch variation & Shift in nominal intensity and defect texture & Lot-stratified validation, recalibration, OOD alarm\\
Model/rule mismatch & Silent false pass after PCB revision & Versioned CAD mask, software/configuration compatibility checks\\
Dataset leakage & Inflated laboratory accuracy & Physical-board and lot keyed splits, audit script\\
Camera/PC failure & Missing inspection or stale result reuse & Trigger sequence IDs, watchdog, fail-safe PLC handshake\\
\bottomrule\end{tabularx}
\end{table}

\subsection{Hardware, economics and commissioning}
An engineering BOM should include camera, lens, tested emission filter, UV source and driver, light shield/interlock, mechanical positioning, PC, I/O, calibration targets, enclosure cooling and traceability storage. For a 70~mm square board, a single 5--12~MP camera and compact directional UV sources may be feasible, but the smallest defect size and sidewall inspection requirements must drive the final optical choice. The earlier conceptual hardware-only budget of roughly €1,500--3,500 per station is \emph{an order-of-magnitude planning assumption}, not a verified October 2026 supplier quotation, not a guaranteed cost, and not inclusive of design engineering, safety validation, line integration, commissioning, process revalidation or service. No production takt or defect sensitivity should be priced as achieved before experiments.

Intel CPUs can run common inference backends; 2026 OpenVINO documentation lists supported Intel Core generations, and ONNX Runtime's execution providers provide portable execution options \citep{openvino2026,onnxproviders}. Candidate model input sizes should be selected after measuring physical defect detectability. An optimized CPU inference engine can still lose to a fixed-rule approach if camera latency, feature tiling or database I/O dominates. A cost-performance plot should compare measured defect-level sensitivity at the selected false-reject rate \emph{versus} whole-cycle 95th-percentile time and amortized station cost. Avoid meaningless comparisons between vendor maximum pixels per second and our end-to-end cycle without identical test boards.

\subsection{Standards and deployment validation}
IPC material qualification (IPC-CC-830C), assembled-board acceptance (IPC-A-610J) and process guidance (IPC-HDBK-830A) play different roles \citep{ipccc830c,ipca610j,ipchdbk830a}. A conformal-coating inspection algorithm is not itself a material qualification procedure. The acceptance recipe must come from authorized full-text requirements and product-specific drawings. Optical and mechanical safety engineering should incorporate IEC~62471 and ISO~12100, with actual measurements and formal assessment \citep{iec62471,iso12100}. Attribute-based lot sampling concepts can support ongoing process monitoring (ISO~2859-1:2026), but must not be conflated with proving sensitivity of the automated 100\%-inspection device \citep{iso2859}.

Pilot operation should initially compare AI decisions against an existing approved inspection process without automatically rejecting production. Review discordant cases, investigate systematic manufacturing shifts, and retrain/revalidate only under controlled revision. Production acceptance requires documented optical calibration, traceability, operator training, maintenance, failure recovery and periodic challenge boards with known defects. A claim of production suitability should only follow these steps.

\section{Research gaps, recommended experiments and conclusions}
\subsection{What prior work already covers}
The search establishes that fluorescent conformal coatings on PCBs have already been inspected using UV candidate extraction, CNN verification, transfer learning, YOLO-style methods and Faster R-CNN \citep{lee2021cnn,lee2021transfer,choi2024uv,wang2025bubble,patel2026coverage}. More general PCB inspection has long used registration, reference difference and learned segmentation \citep{hara1988fluorescent,qub2025registration,ebayyeh2020aoi}. Commercial CCI systems already use coating plans, UV lighting and optional side cameras or thickness metrology \citep{goepelcci,nordsoncci,mirteccci}. These facts exclude simplistic novelty claims that AI, UV, reference masks, image tiling or hybrid pipelines are invented here.

\subsection{Research gaps defensible from the reviewed evidence}
\textbf{Gap 1: reproducible multi-defect evaluation.} The located direct studies generally target selected bubbles or coverage-related defect categories; comparable multi-class, board- and lot-disjoint error rates are not established by their public abstracts.\par
\textbf{Gap 2: observable scope and acceptance semantics.} The manufacturer's CAD-driven approach shows that a coating plan matters; an academic study can quantify the incremental value of explicit required/forbidden/occluded masks at matched false acceptance and rejection.\par
\textbf{Gap 3: illumination robustness under controlled cost.} UV intensity selection has already been studied \citep{choi2024uv}. What remains to test is whether reproducible correction and minimal additional illumination exposures reduce misses enough to offset line-time and hardware complexity.\par
\textbf{Gap 4: CPU edge validation.} Generic AI studies claim high GPU speed, but published results do not establish our complete cycle time on a cheap production-class Intel PC.\par
\textbf{Gap 5: source traceability and transferable datasets.} Existing open datasets cover bare PCB defects and general anomalies rather than a confirmed public, image-level UV coating benchmark \citep{lv2024dspcb,bergmann2019mvtec,visa2022dataset}. A properly documented new dataset is justified; publication rights are a separate requirement.\par

These are \emph{hypothesized research opportunities} relative to the currently reviewed sources, not proof that no other group has addressed them. Full-text searches of licensed proceedings and patents should precede any definitive novelty declaration.

\subsection{Recommended technical build sequence}
\begin{enumerate}[leftmargin=*]
\item \textbf{Optical feasibility bench.} Test the actual coating under 365~nm UV, verify the emission band with appropriate blocking, and quantify field uniformity using one camera. Acquire coated/uncoated references and carefully validated physical defects.
\item \textbf{Classical baseline first.} Implement image registration and flat-field correction, thresholding, morphological region metrics and the approved CAD coverage map. This provides immediate industrial value and an interpretable research benchmark.
\item \textbf{Representative dataset.} Expand to distinct physical boards and batches; label missing coating, overspray/keep-out violations, bubbles and dewetting only where experts can verify each category. Record the uninspectable surfaces.
\item \textbf{Lightweight AI comparison.} Train a small U-Net and one transformer alternative only after defining the baseline failures. Focus on rare, irregular or weakly fluorescent defects, not cosmetic improvement in already simple regions.
\item \textbf{Controlled hybrid comparison.} Test learned masking with and without reference constraints, correction and optional additional exposure. Report causal ablation results and miss cases.
\item \textbf{Edge and industrial qualification.} Export the chosen model to ONNX/OpenVINO, test multiple low-cost CPU systems, measure all timing stages and qualify the procedure against approved production acceptance thresholds.
\end{enumerate}

\subsection{Publication strategy}
A defensible scientific title would be \emph{``Reference-Aware UV Fluorescence Inspection of Conformal-Coated PCBs: A Controlled Comparison of Classical Vision and Lightweight Segmentation''}. Other viable emphases are low-data coating segmentation, multi-angle optical robustness or cost--sensitivity trade-offs. The most compelling paper should present a \emph{specific measured improvement} rather than an untested multi-component system description. Required figures include physical UV enclosure, emission/filter response, calibration map, defect taxonomy with expert labels, architecture, ablation performance curves, size-stratified misses and CPU cycle-time distributions. A release-approved evaluation subset and configuration-controlled source code would improve reproducibility; confidential production images need not be published to establish the methodology.

\subsection{Final synthesis}
For a 70~mm square assembly, controlled UV fluorescence and a single industrial camera offer a plausible economical path to optical verification of \emph{visible} coating coverage. The literature strongly supports building a reference-aware classical baseline before investing in larger AI models. It also demonstrates that the most dangerous research error would be to claim coating thickness, rare-defect reliability or production-ready speed based solely on a bright image or a generic benchmark. The contribution worth testing is an inspection \emph{system of measured optical signals, validated regions, independently verified defects, interpretable decisions and repeatable end-to-end timing}. Until those experiments exist, the present document should be used as a rigorous methodology and informed literature foundation, not as a results paper.

\appendix
\section{Supplementary research instruments and source audit}
\subsection{Companion structured files}
The following CSVs accompany this manuscript in the branch: \texttt{sources/source\_catalog.csv} (verified source identities, URLs, group, access status and cautionary notes); \texttt{sources/search\_log.csv} (query families and date); \texttt{sources/evidence\_matrix.csv} (critical methodological fields for detailed full-text extraction); \texttt{sources/dataset\_register\_template.csv} (prospective specimens/lot and acquisition metadata); and \texttt{sources/experiment\_plan.csv} (planned comparisons, metrics and evidence gates). The citation key is the stable join identifier. A source has not been \emph{fully evaluated} merely because its DOI was found.

\subsection{Source access and interpretation notes}
\begin{enumerate}[leftmargin=*]
\item A publisher page showing an abstract authenticates the publication and abstract-level claims, but does not reveal all acquisition details or test-set composition.
\item The freely exposed pages of IPC-HDBK-830A, IPC-A-610J and IPC-CC-830C are mainly covers/tables of contents; full licensed copies are required to implement exact acceptance thresholds.
\item Manufacturer claims about smallest detectable bubbles and throughput are product specifications, not independent scientific sensitivity with confidence intervals.
\item General PCB copper defect datasets are useful for algorithm prototyping, but not for evidence of conformal-coating performance.
\item Publications about coatings on nozzles and OCT thickness maps are included deliberately as \emph{transfer or complementary modality} evidence.
\item Publication year in bibliographic records is the official cited publication year. Proceedings timing and early-online dates may differ.
\end{enumerate}

\subsection{Search-log outline}
The initial search date, topic-level queries and source categories have been retained in a machine-readable search log. A fully systematic update must export raw database results, count deduplicated records, archive inclusion/exclusion decisions and verify the actual complete text for key studies before paper submission \citep{page2021prisma}.

\subsection{Technical acceptance checklist}
\begin{enumerate}[leftmargin=*]
\item UV leakage and electrical hazards independently assessed; access interlock verified.
\item Full acceptance drawing available and interpreted by qualified manufacturing personnel.
\item Calibration reproducibility and camera/illumination quality gates specified.
\item Ground-truth annotations independently reviewed and physical IDs traceable.
\item Test boards, production lots and image augmentations segregated without leakage.
\item Relevant failure modes have measured recall and uncertainty, not only headline accuracy.
\item Model thresholds selected before inspection of independent test set.
\item Full CPU inspection cycle and fault recovery experimentally measured.
\item Scope of hidden/occluded surfaces disclosed and a manual/alternate inspection method defined.
\item Data confidentiality, software licence and any published dataset release approved.
\end{enumerate}
\subsection{Auditable bibliography inventory}
The companion catalog contains 58 source records with direct links and evidence statuses. The table lists identity and relevant grouping; see the CSV for caveats and URLs.
\begin{longtable}{p{0.10\textwidth}p{0.66\textwidth}p{0.14\textwidth}}
\caption{Initial source inventory, verified at bibliographic-record level on 23 October 2026.}\\
\toprule Key & Publication / official resource & Category \\ \midrule
\endfirsthead
\toprule Key & Publication / official resource & Category \\ \midrule
\endhead
\texttt{lee2021cnn} & A Problematic Bubble Detection Algorithm for Conformal Coated PCB Using Convolutional Neural Networks (2021) & Coating \\
\texttt{lee2021transfer} & A Bubble Detection Method for Conformal Coated PCB Using Transfer Learning based CNN (2021) & Coating \\
\texttt{zouhri2021faster} & Faster R-CNN assessment for air bubbles detection in the conformal coating application (2021) & Coating \\
\texttt{zouhri2024ai} & AI-based Bubbles Detection in the Conformal Coating for Enhanced Quality Control in Electronics Manufacturing (2024) & Coating \\
\texttt{choi2024uv} & Optimal UV Light Intensity Selection Algorithm for Detecting the Tiny Bubble in Conformal Coating (2024) & Coating \\
\texttt{wang2025bubble} & An Advanced YOLOv8-Nano Model with Attention Neck Network for PCB Bubble Detection (2025) & Coating \\
\texttt{patel2026coverage} & Automated Conformal Coating Coverage Defect Detection in Smart Manufacturing Using Faster R-CNN (2026) & Coating \\
\texttt{pidurkar2026hybrid} & A Hybrid Deep-Learning and Image-Processing Approach for Automated PCB Fault Detection (2026) & PCB \\
\texttt{shao2019oct} & Nondestructive Measurement of Conformal Coating Thickness on Printed Circuit Board With Ultra-High Resolution Optical Coherence Tom... (2019) & Metrology \\
\texttt{lf2023oct} & Rapid imaging and product screening with low-cost line-field Fourier domain optical coherence tomography (2023) & Metrology \\
\texttt{han2025coating} & Visual coating inspection framework via self-labeling and multi-stage deep learning strategies (2025) & Transfer \\
\texttt{zi2026coating} & Coating defect detection in intelligent manufacturing: Advances, challenges, and future trends (2026) & Transfer \\
\texttt{hara1988fluorescent} & A System for PCB Automated Inspection Using Fluorescent Light (1988) & Transfer \\
\texttt{crowther2022filters} & Ultraviolet Fluorescence Photography---Choosing the Correct Filters for Imaging (2022) & Optics \\
\texttt{zeissfluorescence} & Principles of Fluorescence Imaging (2019) & Optics \\
\texttt{viscomcci} & Conformal Coating Inspection Provides Greater Safety (2020) & Vendor \\
\texttt{goepelcci} & Conformal Coating Inspection: Multi Line CCI (2026) & Vendor \\
\texttt{nordsoncci} & FX-940UV Automated Conformal Coat Inspection (2026) & Vendor \\
\texttt{mirteccci} & GENESYS-CC Conformal Coating AOI (2026) & Vendor \\
\texttt{ipchdbk830a} & IPC-HDBK-830A: Guidelines for Design, Selection and Application of Conformal Coatings (2013) & IPC \\
\texttt{ipca610j} & IPC-A-610J: Acceptability of Electronic Assemblies (2024) & IPC \\
\texttt{ipccc830c} & IPC-CC-830C: Qualification and Performance of Electrical Insulating Compound for Printed Wiring Assemblies (2019) & IPC \\
\texttt{ipc2024release} & IPC Releases J Revisions to Two Leading Standards for Electronics Assembly (2024) & IPC \\
\texttt{ebayyeh2020aoi} & A Review and Analysis of Automatic Optical Inspection and Quality Monitoring Methods in Electronics Industry (2020) & Review \\
\texttt{zhou2023pcb} & Review of vision-based defect detection research and its perspectives for printed circuit board (2023) & Review \\
\texttt{ieee2023pcb} & A Comprehensive Review of Deep Learning-Based PCB Defect Detection (2023) & Review \\
\texttt{qub2025registration} & Deep learning-enhanced defects detection for printed circuit boards (2025) & PCB \\
\texttt{guo2025cmunet} & CM-UNetv2: An Enhanced Semantic Segmentation Model for Precise PCB Defect Detection and Boundary Restoration (2025) & PCB \\
\texttt{u2d2pcb2024} & U2D2PCB: Uncertainty-Aware Unsupervised Defect Detection on PCB Images Using Reconstructive and Discriminative Models (2024) & PCB \\
\texttt{lim2023pcb} & A deep context learning based PCB defect detection model with anomalous trend alarming system (2023) & PCB \\
\texttt{tang2019deeppcb} & Online PCB Defect Detector on a New PCB Defect Dataset (2019) & Data \\
\texttt{lv2024dspcb} & A dataset for deep learning based detection of printed circuit board surface defect (2024) & Data \\
\texttt{dspcbfigshare} & DsPCBSD+ dataset release (2024) & Data \\
\texttt{jessurun2022fpic} & FPIC: A Novel Semantic Dataset for Optical PCB Assurance (2022) & Data \\
\texttt{bergmann2019mvtec} & MVTec AD -- A Comprehensive Real-World Dataset for Unsupervised Anomaly Detection (2019) & Data \\
\texttt{mvtecdataset} & MVTec AD Dataset and License (2019) & Data \\
\texttt{bergmann2022loco} & Beyond Dents and Scratches: Logical Constraints in Unsupervised Anomaly Detection and Localization (2022) & Data \\
\texttt{visa2022dataset} & Visual Anomaly (VisA) Dataset (2022) & Data \\
\texttt{otsu1979} & A Threshold Selection Method from Gray-Level Histograms (1979) & Method \\
\texttt{fischler1981ransac} & Random Sample Consensus: A Paradigm for Model Fitting with Applications to Image Analysis and Automated Cartography (1981) & Method \\
\texttt{zhang2000calibration} & A Flexible New Technique for Camera Calibration (2000) & Method \\
\texttt{zitova2003registration} & Image registration methods: a survey (2003) & Method \\
\texttt{ronneberger2015unet} & U-Net: Convolutional Networks for Biomedical Image Segmentation (2015) & Method \\
\texttt{chen2018deeplab} & Encoder-Decoder with Atrous Separable Convolution for Semantic Image Segmentation (2018) & Method \\
\texttt{liu2021swin} & Swin Transformer: Hierarchical Vision Transformer Using Shifted Windows (2021) & Method \\
\texttt{xie2021segformer} & SegFormer: Simple and Efficient Design for Semantic Segmentation with Transformers (2021) & Method \\
\texttt{cheng2022mask2former} & Masked-Attention Mask Transformer for Universal Image Segmentation (2022) & Method \\
\texttt{howard2019mobile} & Searching for MobileNetV3 (2019) & Method \\
\texttt{mehta2022mobilevit} & MobileViT: Light-weight, General-purpose, and Mobile-friendly Vision Transformer (2022) & Method \\
\texttt{roth2022patchcore} & Towards Total Recall in Industrial Anomaly Detection (2022) & Method \\
\texttt{zavrtanik2021draem} & DRAEM - A Discriminatively Trained Reconstruction Embedding for Surface Anomaly Detection (2021) & Method \\
\texttt{batzner2024efficientad} & EfficientAD: Accurate Visual Anomaly Detection at Millisecond-Level Latencies (2024) & Method \\
\texttt{iso12100} & ISO 12100:2010 Safety of machinery--Risk assessment and risk reduction (2010) & Standard \\
\texttt{iec62471} & IEC 62471:2006 Photobiological safety of lamps and lamp systems (2006) & Standard \\
\texttt{page2021prisma} & The PRISMA 2020 statement: an updated guideline for reporting systematic reviews (2021) & Review \\
\texttt{openvino2026} & OpenVINO 2026 System Requirements (2026) & Deploy \\
\texttt{onnxproviders} & ONNX Runtime Execution Providers (2026) & Deploy \\
\texttt{iso2859} & ISO 2859-1:2026 Sampling procedures for inspection by attributes (2026) & Standard \\
\bottomrule
\end{longtable}


\clearpage
\phantomsection\addcontentsline{toc}{section}{References}
\bibliography{research/uv_cci_2026/references}
\end{document}
